\scanpage{4}
% A single diagonal cancellation crosses the typescript and reaches into the handwritten draft below it.
\deleted{\uline{Corollaire} \deleted{\unsure{4.2.|4.3.} quater} Soient $A$ un anneau à spectre irréductible, $\underline{p}$ son unique idéal premier minimal, $B$ une $A$-algèbre de présentation finie $M$ un $B$-module de présentation finie. Supposons que $M_{\underline{p}}$ soit plat\inserted{ sur $A_{\underline{p}}$}, alors il existe $f\handwritten{\in}A-\underline{p}$ tel que $\unsure{\mathbb{M}|M}_f$ soit \uline{libre} sur $A_f$.

\deleted{Supposons d'abord $A$ noethérien, donc $A_{\underline{p}}$ local \uline{artinien}, alors l'hypothèse que $M_{\underline{p}}$ est plat sur $A_{\underline{p}}$} En vertu de \deleted{\unsure{4.1. ter|4.1. bis}}\inserted{Prop.A}, quitte à remplacer $A$ par un $A_f$ \deleted{\unsure{sur}} avec $f\quad A-\underline{p}$ convenable, on peut supposer \unsure{qu'i|qu'il} existe un idéal $J$ \inserted{de type fini} dans $A$ tel que, pour toute $A$-algèbre $A'$, on ait $M\otimes_A A'$ est libre sur $A'$, ou encore plat sur $A'$, sss $A\quad A'$ s'annule sur $J$. Appliquant ceci au cas $A'=A_{\underline{p}}$, on voit que l'on a $JA_{\underline{p}}=0$ i.e. $J_{\underline{p}}=0$. Comme $J$ est de type fini, il s'ensuit qu'il existe $f\quad A-\underline{p}$ tel que $J_f=0$ i.e. $JA_f=0$, ce qui implique que $M_f$ est libre sur $A_f$, cqfd

Je m'aperçois à l'instant d'une amélioration dans une autre direction de 4.1. , dont je vais reprendre l'énoncé:

\uline{Proposition} 4.1. Soient $A$ un anneau, $I$ un idéal \deleted{nilpotent} de $A$, $(A_i)_{i\;I}$ une famille de $A$-algèbres telle que l'intersection des noyaux des homomorphismes $A\quad A_i$ soit réduite à 0, $M$ un $A$-module. O\deleted{n}n suppose $M\otimes_A A/I$ libre sur $A/I$ (resp. plat sur $A/I$, \uline{et} l'ensemble d'indices $i$ fini), et pour tout $i$, $M\otimes_A A_i$ libre (resp. plat) sur $A_i$. Alors \inserted{\handwritten{$\forall n$,}} \inserted{\handwritten{$M\otimes_A(A/I^{n+1})$}} $M$ est \deleted{plat sur \unsure{$A$.|$A$}} libre (resp. plat) sur $A\inserted{\handwritten{/I^{n+1}}}$.

Dans le cas non respé, il suffit de prouver aussi la platitude de $M_{\inserted{\handwritten{n}}}$ sur $A_{\inserted{\handwritten{n}}}$, qui entraine la liberté, grâce à \inserted{\handwritten{$M_0=$}}$M\otimes_A A/I$ libre et $I$ nilpotent.

% Successive starts at a handwritten corollary, crossed by a second, looping stroke.
\handwritten{\uline{Corollaire} \deleted{Supposons $A$ anneau \uline{noethérien}\inserted{\unsure{noeth.}}, $M$ un $A$-module,} \margin{$\Bigg[$}\deleted{$(\unsure{\mathfrak{a}|\mathfrak{U}|I}_i)_{i\in\unsure{I}}$ des idéaux de $A$, \inserted{$(\bigcap\mathfrak{a}_i=0)$} \unsure{Supposons} que $\forall i$, $M\otimes_A A/\mathfrak{a}_i$ soit \unsure{plt|plat} sur $\mathfrak{a}_i$. Alors pour tout entier $n$,} Supposons \deleted{\unsure{que}} qu'on est dans l'un des cas \unsure{suivants}

(i) $I$ nilpotent}}%

\handwritten{(ii) $A$ \unsure{noeth.} et $M$ \unsure{idél\textsuperscript{emt}|idéalement|idéliquement} séparé pour $I$.

Alors \deleted{\unsure{les deux conditions sont}}}
\[
\handwritten{M\ \unsure{\text{plt}|\text{plat}}/A
 \iff
 \begin{aligned}
 M_0&\ \text{plt sur }A_0\\
 M_i&\ \text{plt sur }A_i,\ \forall i
 \end{aligned}}
\]
% archive sheet number: 3
